%Master File:lectures.tex

\lesson{Exploring Biological Data}
\vspace{-2cm}
\begin{center}
 \includegraphics[height=8cm]{lda}
\end{center}

\keywords{I am not giving this lecture. Classification, Regression, Decision Trees, LDA, Feature
  Selection}
%%%%%%%%%%%%%%%%%%%%%%% Next Slide %%%%%%%%%%%%%%%%%%%%%%%
\renewcommand{\Outline}{%
\begin{slide}
\section[1]{Outline}

\begin{minipage}{10cm}
  \begin{enumerate}\squeeze
    \outlineitem{Classification}{class}
    \outlineitem{Finding Relevant Variables}{relevance}
    \outlineitem{Feature Selection}{features}
    \outlineitem{Clustering}{clustering}
   \end{enumerate}
\end{minipage}\hfill
\begin{minipage}{11cm}
  \vfill
  \includegraphics[width=11cm]{linear-discriminant-analysis}
  \vfill
\end{minipage}
\end{slide}
\addtocounter{outlineitem}{1}
}

\setcounter{outlineitem}{1}
%%%%%%%%%%%%%%%%%%%%%%% Next Slide %%%%%%%%%%%%%%%%%%%%%%%
\Outline
\toptarget{firstoutline}
%%%%%%%%%%%%%%%%%%%%%%% Next Slide %%%%%%%%%%%%%%%%%%%%%%%

\begin{slide}
\section{Data Exploration}

\begin{PauseHighLight}
  \begin{itemize}
  \item Machine learning is often used to make predictions\pause
  \item But it can also be used for understanding data\pause
  \item You don't need to be a great expert in machine learning
    \begin{itemize}
    \item Use existing tools (e.g. scikit-learn)\pause
    \item Use common sense!\pauseb
    \end{itemize}
  \end{itemize}
\end{PauseHighLight}

\end{slide}

%%%%%%%%%%%%%%%%%%%%%%% Next Slide %%%%%%%%%%%%%%%%%%%%%%%

\begin{slide}
\section[-1]{Coursework}

\begin{PauseHighLight}
  \begin{itemize}
  \item As I am no longer running an assignment this year I am not
    giving this lecture\pause
  \item I leave the lecture on the web pags as it does provide a high
    level introduction to understanding data which some people might
    find useful\pause
  \item You will not be examined on any of this material\pause
  \end{itemize}
\end{PauseHighLight}


\end{slide}

%%%%%%%%%%%%%%%%%%%%%%% Next Slide %%%%%%%%%%%%%%%%%%%%%%%

\begin{slide}
\section{The Data}

\begin{PauseHighLight}
  \begin{itemize}
  \item You are being asked to look at some cancer data\pause
  \item There are different data sets belonging to different
    cancers\pause
  \item Real data is noisy, messy, incomplete, often rather sparse\pause
  \item You are \emph{not} being asked to come up with a good algorithm
    for determining the type of cancer\pause
  \item You are being asked to explore the data and see what it tells
    you (if anything)\pause
  \end{itemize}
\end{PauseHighLight}

\end{slide}

%%%%%%%%%%%%%%%%%%%%%%% Next Slide %%%%%%%%%%%%%%%%%%%%%%%

\begin{slide}
\section{Classification}

\begin{PauseHighLight}
  \begin{itemize}
  \item One of the simplest analysis is to see if you can make a
    classifier that makes predictions\pause
  \item There are lots of classifies that do this (use someone else's)\pause
    \begin{itemize}\squeeze
    \item KNN
    \item Perceptron, logistic regression, MLP
    \item Decision trees, random forest, XGBoost
    \item SVM, \ldots\pause
    \end{itemize}
  \end{itemize}
\end{PauseHighLight}

\end{slide}


%%%%%%%%%%%%%%%%%%%%%%% Next Slide %%%%%%%%%%%%%%%%%%%%%%%

\begin{slide}
\section[-2]{$K$-Nearest Neighbours}
\pause
\pauselevel{=1}
\pb
\begin{center}
  \multipdf[width=0.6\linewidth]{knn}\pause
\end{center}
\end{slide}

%%%%%%%%%%%%%%%%%%%%%%% Next Slide %%%%%%%%%%%%%%%%%%%%%%%

\begin{slide}
\section[-2]{Preparing Data}
  
\begin{PauseHighLight}
  \begin{itemize}
  \item Before running a learning machine you need to look at the data
    (pandas isn't bad for this)\pause
  \item The data might be numeric or categorical\pause
  \item Often you might have missing data\pause
  \item Many techniques require numeric values (can you turn categorical
    data into numerical data?)\pause
  \item Is your data normalised (are they over roughly same range?)\pause
  \item Missing data can sometimes be replaced (e.g. with mean
    values)\pause
  \item Decision trees can cope both with categorical data and missing
    values\pause 
  \end{itemize}
\end{PauseHighLight}

\end{slide}


%%%%%%%%%%%%%%%%%%%%%%% Next Slide %%%%%%%%%%%%%%%%%%%%%%%

\begin{slide}
\section{Binary or Multi-Class}

\begin{PauseHighLight}
  \begin{itemize}
  \item There are lots of different classes\pause---you don't need to
    look at all of them\pause
  \item You could just try to identify one class, using all the others
    as examples of not that class (SVMs, perceptrons and logistic
    regression do binary classification)\pause
  \item You can look at multi-class trying to classify the different
    classes (you might want to use KNN, a MLP with \textit{softmax} or
    decision tree based algorithms)\pause
  \end{itemize}
\end{PauseHighLight}

\end{slide}

%%%%%%%%%%%%%%%%%%%%%%% Next Slide %%%%%%%%%%%%%%%%%%%%%%%

\begin{slide}
\section{Balancing Datasets}

\begin{PauseHighLight}
  \begin{itemize}
  \item Often learning machines struggle doing classification when the
    datasets are very unbalanced\pause
  \item In binary classification you might have many more positive
    examples than negative examples\pause
  \item In multi-class classification you might have some classes with
    many more examples than others\pause
  \item In this case it can sometimes help to balance the datasets
    either
    \begin{itemize}
    \item By throwing away some data or
    \item Using data of rare classes multiple times\pause
    \end{itemize}
  \end{itemize}
\end{PauseHighLight}

\end{slide}

%%%%%%%%%%%%%%%%%%%%%%% Next Slide %%%%%%%%%%%%%%%%%%%%%%%

\begin{slide}
\section{Measuring Prediction Performance}

\begin{PauseHighLight}
  \begin{itemize}
  \item To see whether we can predict a cancer we must test our
    classifier on unseen data\pause
  \item To do this we have to split the data into a \emph{training} and
    \emph{test} set\pause
  \item We use the training set to train the data and the test set to
    test the algorithm\pause
  \item There is a dilemma: when we have a small amount of data
    \begin{itemize}\squeeze
    \item We want to use as much data as possible for learning (to get
      good performance)
    \item We want to use as much data as possible for testing (to get an
      accurate measure of performance)\pause
    \end{itemize}
  \end{itemize}
\end{PauseHighLight}

\end{slide}


%%%%%%%%%%%%%%%%%%%%%%% Next Slide %%%%%%%%%%%%%%%%%%%%%%%

\begin{slide}
\section[-2]{Cross Validation}

\pb
\begin{itemize}
\item We can solve our dilemma by repeating the training many times on
  different partitionings\pauseh
\item This is known as $K$-fold cross validation\pauseh
\end{itemize}
\begin{center}
  \includegraphics[width=0.9\linewidth]{cross_validation0}\mypl{1 :3}
  \multido{\ia=1+1,\ib=4+1}{23}{%
    \llap{\includegraphics[width=0.9\linewidth]{cross_validation\ia}}\mypl{\ib}}
  \llap{\includegraphics[width=0.9\linewidth]{cross_validation23}}
\end{center}
\begin{itemize}
\item Leave-one-out cross-validation is extreme case\pause
\end{itemize}
\end{slide}

%%%%%%%%%%%%%%%%%%%%%%% Next Slide %%%%%%%%%%%%%%%%%%%%%%%

\begin{slide}
\section{Price of Cross Validation}

\begin{PauseHighLight}
  \begin{itemize}
  \item Cross-validation is computationally expensive ($K$-fold
    cross-validation requires $K$ times as much work)\pause
  \item There are attempts at estimating generalisation error more
    cheaply (boot-strapping) methods, but these are not very
    accurate\pause
  \item Cross-validation is only necessary when you have little
    data\pause 
  \item Re-running code is usually cheap compared with writing the
    code\pause
  \end{itemize}
\end{PauseHighLight}

\end{slide}

%%%%%%%%%%%%%%%%%%%%%%% Next Slide %%%%%%%%%%%%%%%%%%%%%%%
\Outline %% Finding Relevant Variables
%%%%%%%%%%%%%%%%%%%%%%% Next Slide %%%%%%%%%%%%%%%%%%%%%%%

\begin{slide}
\section{What's Relevant}

\begin{PauseHighLight}
  \begin{itemize}
  \item Not all variables are necessarily relevant\pause
  \item There are a couple of quick ways of seeing which variables are
    most relevant
    \begin{itemize}
    \item Logistic Regression
    \item Decision Trees\pause
    \end{itemize}
  \end{itemize}
\end{PauseHighLight}

\end{slide}

%%%%%%%%%%%%%%%%%%%%%%% Next Slide %%%%%%%%%%%%%%%%%%%%%%%

\begin{slide}
\section{Logistic Regression}

\begin{PauseHighLight}
  \begin{itemize}
  \item In logistic regression we learn a function
    \begin{align*}
      f(\bm{x}) &= \mathrm{sigmoid}\!\left(\bm{w}^\tr\bm{x} +b\right)
      &
      \mathrm{sigmoid}(z) &= \frac{1}{1 + \e{-z}}  \pause
    \end{align*}
  \item The variables, $x_i$, with the largest (absolute) weights,
    $|w_i|$, are likely to be the most important\pause{} (provided you
    have normalised your variables)\pauseb
  \item This is a very standard algorithm that exists in most package,
    e.g.{} scikit-learn (in python)\pause
  \end{itemize}
\end{PauseHighLight}

\end{slide}

%%%%%%%%%%%%%%%%%%%%%%% Next Slide %%%%%%%%%%%%%%%%%%%%%%%

\begin{slide}
\section{Decision Trees}

\begin{PauseHighLight}
  \begin{itemize}
  \item Decision trees builds a binary tree to partition the data,
    $\mathcal{D}=\{(\bm{x}, y)\}$, into the leaves of the tree\pause
  \item Each decision rule depends on a single feature\pause
  \item At each step the rule is chosen that maximise the
    ``\textit{purity}'' of the leaf nodes\pause
  \item Decisions can be made on numerical values or categories\pause
  \end{itemize}
\end{PauseHighLight}

\end{slide}

%%%%%%%%%%%%%%%%%%%%%%% Next Slide %%%%%%%%%%%%%%%%%%%%%%%

\begin{slide}
\section[-2]{Partitioning}

\begin{PauseHighLight}
  \begin{itemize}
  \item Consider a classification problems with examples $(\bm{x}, y)$
    belonging to some classes $y\in\mathcal{C}$\pause
  \item The data is partitioned by the tree into leaves
    \begin{center}
      \includegraphics[width=0.62\linewidth]{treePicture}\pause
    \end{center}
    \vspace*{-1.6em}

  \item The proportion of data points in leaf $\mathcal{L}$ belonging to
    class $c$ is
    \begin{align*}
      p_{c}(\mathcal{L}) = \frac{1}{|\mathcal{L}|}
      \sum_{(\bm{x},y)\in\mathcal{L}} \pred{y = c}
    \end{align*}
    where $\pred{y = c} = 1$ if $y=c$ and 0 otherwise\pause
  \end{itemize}
\end{PauseHighLight}

\end{slide}

%%%%%%%%%%%%%%%%%%%%%%% Next Slide %%%%%%%%%%%%%%%%%%%%%%%

\begin{slide}
\section{Leaf Purity}

\begin{PauseHighLight}
  \begin{itemize}
  \item Two different purity measures, $Q_m(\mathcal{L})$, for a leaf node
    $\mathcal{L}$ are commonly used\pause
    \begin{itemize}
    \item \emph{Gini index}
      \begin{rightImage}{gini}
        \begin{align*}
          Q^g_m(\mathcal{L}) = \sum_{c\in\mathcal{C}} p_{c}(\mathcal{L})\,
          \left(1-p_{c}(\mathcal{L})\right)\pause
        \end{align*}
      \end{rightImage}
    \item \emph{Cross-entropy}
      \begin{rightImage}{crossEnt}
        \begin{align*}
          Q^e_m(\mathcal{L}) = - \sum_{c\in\mathcal{C}} p_{c}(\mathcal{L})\,
          \logg{p_{c}(\mathcal{L})} \pause        
        \end{align*}
      \end{rightImage}
    \end{itemize}
  \end{itemize}
\end{PauseHighLight}

\end{slide}


%%%%%%%%%%%%%%%%%%%%%%% Next Slide %%%%%%%%%%%%%%%%%%%%%%%

\begin{slide}
\section[-1]{Building Decision Trees}

\pb\pause\pauselevel{=1}
\begin{center}
  \multipdf[width=\linewidth]{decisionTree}\pause
\end{center}

\end{slide}

%%%%%%%%%%%%%%%%%%%%%%% Next Slide %%%%%%%%%%%%%%%%%%%%%%%

\begin{slide}
\section{Observations}

\begin{PauseHighLight}
  \begin{itemize}
  \item Decision trees are very useful for exploring new data
    sets\pause---the tree shows what features are most
    important\pause
  \item That is, features near the root of the tree are most
    informative\pause
  \item CART is a classic implementation that builds
    \emph{C}lassification \emph{A}nd \emph{R}egression
    \emph{T}rees\pause
  \item Decision trees handle missing data well\pause
  \item A single decision is usually not that good, but an ensemble of
    decision trees (random forest or boosting) can be very good\pause
  \end{itemize}
\end{PauseHighLight}

\end{slide}

%%%%%%%%%%%%%%%%%%%%%%% Next Slide %%%%%%%%%%%%%%%%%%%%%%%
\Outline %% Feature Selection
%%%%%%%%%%%%%%%%%%%%%%% Next Slide %%%%%%%%%%%%%%%%%%%%%%%

\begin{slide}
\section{Feature Selection}

\begin{PauseHighLight}
  \begin{itemize}
  \item In biology many features may be irrelevant to
    classification\pause
  \item You could try every combination of feature, but for $p$ features
    there are $2^p$ combinations\pause
  \item There are various schemes for trying to find which features are
    important\pause
  \end{itemize}
\end{PauseHighLight}

\end{slide}

%%%%%%%%%%%%%%%%%%%%%%% Next Slide %%%%%%%%%%%%%%%%%%%%%%%

\begin{slide}
\section{Correlation with Class}

\begin{PauseHighLight}
  \begin{itemize}
  \item You can measure the correlation between each feature and the
    class category (assuming binary classification $y_i\in\{-1,1\}$)
    \begin{align*}
      r = \frac{\av{ (x_i-\av{x_i}) (y_i - \av{y_i})
      }}{\sqrt{\Var{x_i}\,\Var{y_i}}}\pause
    \end{align*}
  \item But be aware
    \begin{itemize}
    \item If two variables are correlated with each other both may not
      provide any more information than one alone\pause
    \item Two variables that aren't correlated with the data, might
      together predict classes perfectly\pause
    \end{itemize}
  \end{itemize}
\end{PauseHighLight}

\end{slide}

%%%%%%%%%%%%%%%%%%%%%%% Next Slide %%%%%%%%%%%%%%%%%%%%%%%

\begin{slide}
\section{Building Classifiers}

\begin{PauseHighLight}
  \begin{itemize}
  \item We can systematically build a classifier starting from the most
    informative variable\pause
  \item We then add all possible variables and choose the best to
    form a two variable classifier\pause
  \item This can be repeated until we find that adding additional
    variables doesn't improve generalisation performance\pause
  \item We can also start with all variables and eliminate the worst
    variables one by one\pause
  \end{itemize}
\end{PauseHighLight}

\end{slide}

%%%%%%%%%%%%%%%%%%%%%%% Next Slide %%%%%%%%%%%%%%%%%%%%%%%

\begin{slide}
\section{Feature Selection}

\begin{PauseHighLight}
  \begin{itemize}
  \item There are whole books written on feature selection\pause
  \item Techniques exist, such as Lasso (adding an $L_1$ regulariser),
    that are designed to do automatic feature selection\pause
  \item Often in biology we want to know if we can make reasonably
    reliable predictions with just a few variables\pause
  \item Don't necessarily need very sophisticated techniques to do
    this\pause
  \end{itemize}
\end{PauseHighLight}

\end{slide}


%%%%%%%%%%%%%%%%%%%%%%% Next Slide %%%%%%%%%%%%%%%%%%%%%%%
\Outline %% LDA
%%%%%%%%%%%%%%%%%%%%%%% Next Slide %%%%%%%%%%%%%%%%%%%%%%%

\begin{slide}
\section{Clustering}

\begin{PauseHighLight}
  \begin{itemize}
  \item In determining whether a cell is cancerous it is often useful to
    find projections of the data where clustering is more obvious\pause
  \item A projection is a set of ortho-normal vectors
    \begin{align*}
      \mat{P} =
      \begin{pmatrix}
        \bm{p}_1 & \bm{p}_2 & \cdots & \bm{p}_k
      \end{pmatrix}
    \end{align*}
    such that we can project a $p$-dimensional vector, $\bm{x}$, to a
    $k$-dimensional vector, $\bm{z} = \mat{P}^\tr \bm{x}$\pause
  \end{itemize}
\end{PauseHighLight}

\end{slide}

%%%%%%%%%%%%%%%%%%%%%%% Next Slide %%%%%%%%%%%%%%%%%%%%%%%

\begin{slide}
\section{Fisher Discriminant Analysis \& LDA}

\begin{PauseHighLight}
  \begin{itemize}
  \item If we want to find the best 1-d projection that maximises the
    ratio of inter-class distance to intra-class distances then we use
    Fisher Discriminant Analysis\pause
  \item We can generalise this to the best $k$-d projection using
    \textit{linear discriminant analysis} or LDA\pause
  \item We can use LDA for multiple classes\pause
  \item There are standard packages for running LDA (beware not to
    confuse linear discriminant analysis with latent Dirichlet
    allocation)\pause
  \end{itemize}
\end{PauseHighLight}

\end{slide}

%%%%%%%%%%%%%%%%%%%%%%% Next Slide %%%%%%%%%%%%%%%%%%%%%%%

\begin{slide}
\section{Linear Discrimant Analysis}
\begin{center}
 \includegraphics[width=\linewidth]{lda}
\end{center}

\end{slide}

%%%%%%%%%%%%%%%%%%%%%%% Next Slide %%%%%%%%%%%%%%%%%%%%%%%

\begin{slide}
\section{Visualisation}

\begin{PauseHighLight}
  \begin{itemize}
  \item Using LDA in 2 or 3 dimensions often allows you to visualise the
    clustering of data\pause
  \item This can give an idea of what cancers are like each other\pause
  \item LDA mixes different variables so doesn't help with feature
    selection\pause
  \item Other techniques such as \emph{t-SNE} are also useful for
    visualising cluster\pause
  \end{itemize}
\end{PauseHighLight}

\end{slide}

%%%%%%%%%%%%%%%%%%%%%%% Next Slide %%%%%%%%%%%%%%%%%%%%%%%

\begin{slide}
\section{Conclusions}

\begin{PauseHighLight}
  \begin{itemize}
  \item We've covered some techniques for understand data\pause
  \item There are many more you could try\pause
  \item But you don't need to try everything\pause
  \item Write programs (scripts) for doing data analysis as you will
    need the over and over again\pause
  \item Visualise as much as possible\pause
  \item Python can be very good to do this\pause
  \end{itemize}
\end{PauseHighLight}

\end{slide}


